\subsection{Extended Bayesian information criterion}

The Gaussian prior of $\beta$ is a conjugate prior in \eqref{eq:spline_regression}, yielding a closed expression of $p(y|k,\xi)$. Generally, $p(y|k,\xi)$ is analytically intractable except for the normal regression model. %Then we have to calculate the marginal likelihood numerically. 
For those non-conjugate cases, we utilize the extended Bayesian information criterion (EBIC) of \citet{10.1093/biomet/asn034} to approximate the posterior density. %Informally, the EBIC takes into account the dimension of unknown parameters and the complexity of the model space, extremely useful for variable selection in problems with a huge number of candidate covariates. %\citep{NIPS2010_072b030b, bb6df906-631b-39ae-98e5-1c3a3cd1049c, Luo2015}.
In the spline knot estimation, the cardinality of all candidate knots (\emph{i.e.}, $n$) can be very large but the number of the true knots (\emph{i.e.}, $k$) is small compared to the sample size (\emph{i.e.}, $m$). Thus, EBIC will be extremely useful for model selection as the small-$m$-large-$n$ assumption holds and the Laplace approximation is valid \citep{NIPS2010_072b030b, bb6df906-631b-39ae-98e5-1c3a3cd1049c, Luo2015}.
% give an excellent estimation of the posterior density.

According to the definition, the EBIC of $k,\xi$ is
\begin{equation}
    \label{eq:ebic}
    \text{BIC}_{\gamma}(k,\xi) = -2\log L(\hat{\beta},\hat{\sigma}|y,k,\xi) + (\nu+1) \log m + 2\gamma \log \tau(\mathcal{M}_k),\quad 0\leq \gamma \leq 1,
\end{equation}
where $\hat{\beta},\hat{\sigma}$ are the maximum likelihood estimators of $\beta,\sigma$ given $k,\xi$. Especially in \eqref{eq:spline_regression}, $\hat{\beta}=(Z^\top Z)^{-1}Z^\top y$ and $\hat{\sigma}^2 = y^\top (I_m-Z(Z^\top Z)^{-1}Z^\top)y/m$. %And $\gamma\in [0,1]$ is a hyper-parameter of the penalty term involved with the complexity of $\mathcal{M}_k$. 
As $\gamma=0$, \eqref{eq:ebic} is the ordinary BIC. Since $n_1,\ldots,n_d>>m$, the EBIC with $\gamma>0$ is preferable to the ordinary BIC in knot estimation. From the Laplace approximation, $p(k,\xi|y)\approx \exp\{-\text{BIC}_\gamma(k,\xi)/2\}$, denoted by $\hat{p}(k,\xi|y)$.
%Furthermore, $\text{BIC}_\gamma(k,\xi)$ is the Laplace approximation of $-2\log p(k,\xi|y)$. Thus, $p(k,\xi|y)\approx \exp\{-\text{BIC}_\gamma(k,\xi)/2\}$, denoted by $\hat{p}(k,\xi|y)$.

\begin{lemma}
    In multivariate spline model \eqref{eq:spline_regression}, the EBIC approximation of the posterior density is
    \begin{equation}
        \label{eq:approx}
        \hat{p}(k,\xi|y)\propto m^{-(\nu+1)/2} (\hat{\sigma}^2)^{-m/2} \tau(\mathcal{M}_k)^{-\gamma}.
    \end{equation}
\end{lemma}
Suppose $k',\xi'$ are another group of knots. Comparing \eqref{eq:posterior} and \eqref{eq:approx}, we can find $p(k,\xi|y)/p(k',\xi'|y)\approx \hat{p}(k,\xi|y)/\hat{p}(k',\xi'|y)$ when $m$ is sufficiently large. Since the sampling procedure of MCMC is determined by the posterior density ratio, the EBIC contributes to a consistent estimation.